\documentclass[10pt,a4paper]{article}

\usepackage[margin=1.1cm]{geometry}
\usepackage{fontspec}
\usepackage{xcolor}
\usepackage{enumitem}
\usepackage{hyperref}
\usepackage{microtype}

\IfFileExists{/usr/share/fonts/liberation-sans/LiberationSans-Regular.ttf}{
  \setmainfont[
    Path=/usr/share/fonts/liberation-sans/,
    BoldFont=LiberationSans-Bold.ttf,
    ItalicFont=LiberationSans-Italic.ttf,
    BoldItalicFont=LiberationSans-BoldItalic.ttf
  ]{LiberationSans-Regular.ttf}
}{
  \setmainfont{TeX Gyre Heros}
}
\definecolor{accent}{HTML}{4F46E5}
\definecolor{muted}{HTML}{475569}
\hypersetup{
  colorlinks=true,
  urlcolor=accent,
  linkcolor=accent,
  pdfauthor={Yu Zhu},
  pdftitle={Yu Zhu - Curriculum Vitae}
}

\pagestyle{empty}
\setlength{\parindent}{0pt}
\setlength{\parskip}{1pt}
\setlist[itemize]{leftmargin=1.25em,itemsep=0pt,topsep=1pt,parsep=0pt}
\makeatletter
\renewcommand\section{\@startsection{section}{1}{0pt}%
  {4pt}{2pt}{\large\bfseries\color{accent}}}
\makeatother

\newcommand{\entry}[4]{%
  \textbf{#1}\hfill\textcolor{muted}{#2}\\
  \textit{#3}\hfill\textcolor{muted}{#4}\par
}

\begin{document}

\begin{center}
  {\LARGE\bfseries Yu Zhu}\\[3pt]
  \textcolor{muted}{AI for Science · Computational Chemistry · Computer Vision}\\[4pt]
  \href{mailto:v.yuzhu1@gmail.com}{v.yuzhu1@gmail.com}
  \quad|\quad
  \href{https://github.com/ZeHeru}{github.com/ZeHeru}
  \quad|\quad
  \href{https://scholar.google.com/citations?user=GtyV_YsAAAAJ&hl=en}{Google Scholar}
  \quad|\quad
  \href{https://orcid.org/0009-0007-1068-4548}{ORCID}
\end{center}
\small

\section{Research Profile}
Research at the intersection of AI for Science, computational chemistry, and
computer vision, spanning machine learning force fields, physics-driven
reconstruction of molecular orbitals from scanning tunneling microscopy
images, open-set fine-grained recognition, and materials modeling.

\section{Education}
\entry{M.Sc., Department of Chemistry}{Sep 2023 -- Jun 2026}
      {Fudan University}{Shanghai, China}
\begin{itemize}
  \item Supervisors: Prof. Xin Xu and Young Researcher Sai Duan.
  \item Research: machine learning force fields, open-set recognition, and scanning tunneling microscopy.
\end{itemize}

\entry{B.Sc., College of Chemistry, Chemical Engineering and Materials Science}{Sep 2019 -- Jun 2023}
      {Soochow University}{Suzhou, China}
\begin{itemize}
  \item GPA: 3.8/4.0 (92.1/100).
  \item Supervisors: Prof. Jianlin Yao and Assoc. Prof. Minmin Xu.
\end{itemize}

\section{Experience}
\entry{AI for Science Intern}{Jul 2024 -- Jun 2025}
      {Microsoft Research}{Shanghai, China}
\begin{itemize}
  \item Contributed benchmarks and scalable computational workflows for materials modeling.
  \item Evaluated universal machine learning force fields and ran inference on models up to 1B parameters.
  \item Managed high-throughput datasets with up to 200M entries and built Azure pipelines supporting up to 10K concurrent tasks.
\end{itemize}

\section{Selected Projects}
\textbf{\href{https://github.com/ZeHeru/STM_net/tree/main}{STM-Net: Reconstructing Molecular Orbitals from STM Images}}
\hfill\textcolor{muted}{Jan 2024 -- Jul 2024}
\begin{itemize}
  \item Simulated STM images at the DFT level using Bardeen's approximation.
  \item Developed a physics-driven U-Net to separate s-wave molecular-orbital signals from p-wave CO-tip contributions.
  \item Reached PSNR above 37 dB using 10\% of the training data and above 43 dB on the full dataset.
\end{itemize}

\textbf{\href{https://github.com/microsoft/mattersim}{MatterSim: Materials Modeling Workflows}}
\hfill\textcolor{muted}{Jul 2024 -- Jun 2025}
\begin{itemize}
  \item Built benchmarks and high-throughput workflows using VASP, atomate2, phonopy, phono3py, ESEN, MACE, and ORB.
\end{itemize}

\textbf{MPBR for Open-Set Fine-Grained Recognition}
\hfill\textcolor{muted}{Jan 2025 -- Jun 2025}
\begin{itemize}
  \item Developed benchmark pipelines and evaluated OSR models on iNat2021-OSR using t-SNE, Grad-CAM, Swin Transformer, and HAN-OSFGR.
\end{itemize}

\textbf{SERS Imaging for Fingerprint Identification}
\hfill\textcolor{muted}{Apr 2021 -- Jun 2023}
\begin{itemize}
  \item Performed SERS imaging and chemical mapping of latent and live fingerprints on flexible substrates.
\end{itemize}

\section{Publications}
\begin{enumerate}[leftmargin=1.5em,itemsep=3pt,topsep=2pt]
  \item \textbf{Yu Zhu}, R. Xue, H. Ren, Y. Chen, W. Yan, B. Wu, S. Duan,
  H. Zhang, L. Chi, and X. Xu.
  \href{https://doi.org/10.1021/jacsau.5c00310}{Reconstructing Pristine Molecular Orbitals from Scanning Tunneling Microscope Images via Artificial Intelligence Approaches}.
  \textit{JACS Au}, 5(7), 3163--3170, 2025.

  \item J. Li, Z. Chen, Q. Wang, H. Yang, Z. Lu, G. Li, S. Chen,
  \textbf{Yu Zhu}, et al.
  \href{https://doi.org/10.48550/arXiv.2503.11568}{Probing the Limit of Heat Transfer in Inorganic Crystals with Deep Learning}.
  arXiv:2503.11568, 2025.

  \item J. Tan, P. Jing, \textbf{Yu Zhu}, and Y. Liu.
  \href{https://openaccess.thecvf.com/content/ICCV2025/html/Tan_MPBR_Multimodal_Progressive_Bidirectional_Reasoning_for_Open-Set_Fine-Grained_Recognition_ICCV_2025_paper.html}{MPBR: Multimodal Progressive Bidirectional Reasoning for Open-Set Fine-Grained Recognition}.
  \textit{ICCV}, 1282--1291, 2025.

  \item M. Wu, J. Tan, \textbf{Yu Zhu}, and J. Ren.
  \href{https://doi.org/10.48550/arXiv.2604.03547}{KappaFormer: Physics-aware Transformer for Lattice Thermal Conductivity via Cross-domain Transfer Learning}.
  arXiv:2604.03547, 2026.
\end{enumerate}

\section{Honors \& Awards}
\begin{itemize}
  \item China National Scholarship, top 0.2\% (2025 and 2022).
  \item Silver Award, China International College Students' Innovation and Entrepreneurship Competition, top 0.02\% (2021).
  \item First Prize, Chinese Mathematics Competition for College Students, top 5.3\% (2020).
\end{itemize}

\section{Technical Skills}
\textbf{Scientific computing:} Python, VASP, atomate2, joblib, phonopy, phono3py.\\
\textbf{Machine learning:} U-Net, DAE, Mask2Former, ESEN, MACE, ORB, Swin Transformer, HAN-OSFGR.\\
\textbf{Scientific methods:} STM simulation, Bardeen approximation, Tersoff-Hamann theory, Chen derivative rule, SERS mapping.\\
\textbf{Visualization:} t-SNE, Grad-CAM, Origin.

\end{document}
